\documentclass[conference]{IEEEtran}
\IEEEoverridecommandlockouts

\usepackage{cite}
\usepackage{amsmath,amssymb,amsfonts}
\usepackage{graphicx}
\usepackage{textcomp}
\usepackage{xcolor}
\usepackage{url}
\usepackage{array}
\usepackage{multirow}
\usepackage{balance}

\def\BibTeX{{\rm B\kern-.05em{\sc i\kern-.025em b}\kern-.08em
    T\kern-.1667em\lower.7ex\hbox{E}\kern-.125emX}}

\begin{document}

\title{Source-Grounded, Multi-Modal Retrieval-Augmented\\
Question Answering: A Survey of Techniques, Systems,\\
and Security Considerations, with SAGE as a Case Study}

\author{
\IEEEauthorblockN{Srujan Zanjal}
\IEEEauthorblockA{
Department of Computer Engineering\\
St. Vincent Pallotti College of Engineering and Technology\\
Nagpur, India\\
zanjalsrujan94@gmail.com}
}

\maketitle

\begin{abstract}
Retrieval-augmented generation (RAG) has become the dominant practical
strategy for grounding large language model outputs in external
evidence, but the literature that supports it is fragmented across at
least five largely separate threads: core RAG architectures, modality-specific
extraction (documents, speech and video, and source code), multi-source
``notebook-style'' assistants and general-purpose orchestration
frameworks, hallucination and confidence estimation, and -- increasingly
urgent as these systems are given access to arbitrary external content --
the security of the retrieval pipeline itself. This survey reviews each
thread, synthesises the design trade-offs that recur across it, and
argues that a practical multi-modal, source-grounded assistant is best
understood as a set of choices along four axes: how broad, low-specificity
questions are made retrievable; how citation provenance is preserved
per modality; how multiple heterogeneous sources are fairly represented
in a combined answer; and how the ingestion and generation steps are
hardened against content the system does not control. We use SAGE, a
prototype multi-source knowledge engine that ingests websites, PDFs,
YouTube videos, and public GitHub repositories, as a running case study
that instantiates concrete answers to all four axes -- deterministic
source-overview chunks, modality-aware citation rendering, diversity-floor
reranking, and combined SSRF and indirect-prompt-injection defenses --
and compare it qualitatively against a closed consumer assistant
(NotebookLM) and general-purpose open-source RAG frameworks (LangChain,
LlamaIndex, Haystack). We close with open challenges -- calibrated
confidence, hybrid lexical-semantic retrieval, and controlled multi-modal
benchmarks -- that remain unresolved across the surveyed systems.
\end{abstract}

\begin{IEEEkeywords}
retrieval-augmented generation, survey, source-grounded question
answering, multi-modal retrieval, hallucination, prompt injection,
vector databases, document intelligence, video intelligence, code
intelligence
\end{IEEEkeywords}

\section{Introduction}

Large language models answer fluently even when they should not.
Retrieval-augmented generation (RAG) responds to this by conditioning
generation on passages retrieved from an external corpus rather than
relying only on parameters learned at training time, and has been shown
to reduce unsupported claims on knowledge-intensive tasks relative to a
purely parametric model \cite{lewis2020rag}. In the years since, RAG has
moved from a single retrieve-then-generate recipe to a broad family of
architectures, and the corpora it grounds against have expanded well
beyond short text passages to documents with layout, spoken and video
content, and source code repositories. A parallel, more applied
literature has produced consumer and developer-facing systems --
notebook-style multi-source assistants such as NotebookLM
\cite{notebooklm2026}, and general-purpose orchestration frameworks such
as LangChain \cite{langchain2026}, LlamaIndex \cite{llamaindex2026}, and
Haystack \cite{haystack2026} -- that package these architectural ideas
into usable products and libraries.

This survey has two goals. The first is to consolidate five threads of
work that are usually read separately -- RAG foundations, modality-specific
extraction, multi-source assistants and frameworks, hallucination and
confidence estimation, and RAG-specific security -- into a single
account of the design space a builder of a multi-modal, source-grounded
assistant actually faces. The second is to ground that account in a
concrete system. We use \textbf{SAGE}, a prototype we built that ingests
websites, PDFs, YouTube videos, and public GitHub repositories into one
retrieval and citation pipeline, as a case study throughout Section IX,
not as the survey's primary contribution but as a way of making the
trade-offs discussed in Sections III--VIII concrete and checkable
against a working implementation rather than only against claims made
in the surveyed papers themselves.

We organise the recurring design space around four questions that, in
our reading of the literature, separate a system that merely retrieves
passages from one that behaves like a trustworthy source-grounded
assistant: (i) how are broad, low-specificity questions (``what is this
about'') made retrievable at all, given that dense retrieval over
fine-grained chunks tends to under-serve them; (ii) how is citation
provenance preserved in a form appropriate to the source's modality,
since a URL, a page number, a timestamp, and a file/line range are not
interchangeable; (iii) when a query spans multiple sources, how is fair
representation across sources guaranteed rather than left to whichever
source's chunks happen to score highest; and (iv) given that ingestion
fetches content the system does not control, how is the pipeline
defended against that content actively working against the system. We
return to these four questions in Section IX as the specific criteria
against which SAGE, NotebookLM, and generic framework applications are
compared.

The rest of this paper is organised as follows. Section II states the
survey's scope. Section III reviews core RAG architectures. Section IV
covers modality-specific retrieval and extraction for documents, video
and speech, and code. Section V reviews multi-source assistants and
orchestration frameworks. Section VI reviews confidence and hallucination
in RAG. Section VII reviews security considerations specific to
retrieval pipelines. Section VIII reviews hybrid and advanced retrieval
strategies. Section IX presents SAGE as a case study and a comparative
table. Section X discusses open challenges, and Section XI concludes.

\section{Scope and Method}

This survey focuses on retrieval-augmented systems that ground answers
in externally supplied evidence -- as opposed to purely parametric LLM
prompting, in-context few-shot learning without retrieval, or fine-tuning
approaches to knowledge injection, which we treat as out of scope except
where they interact directly with retrieval. Within that focus we
prioritise: (a) foundational architectures that later applied work
consistently builds on; (b) modality coverage that maps to the four
source types a multi-modal assistant like SAGE actually has to
handle -- web pages, documents, video/speech, and code; (c) systems and
frameworks that are either widely deployed (NotebookLM) or widely used
as building blocks (LangChain, LlamaIndex, Haystack); and (d) two
cross-cutting concerns -- faithfulness/hallucination and security -- that
apply regardless of modality and that we found were often treated as
secondary in modality-specific papers despite being first-order concerns
for any deployed system. Coverage is necessarily illustrative rather than
exhaustive given the pace of publication in this area; we favour
frequently-cited foundational work and recent surveys
\cite{gao2023ragsurvey,slr2025rag,ragevalsurvey2025} over an attempt to
enumerate every recent variant.

\section{Foundations of Retrieval-Augmented Generation}

\subsection{From Parametric to Retrieval-Augmented Generation}
RAG's original formulation pairs a pretrained sequence-to-sequence
generator with a differentiable retriever over a document index, trained
so that the generator learns to make use of whatever the retriever
surfaces \cite{lewis2020rag}. Dense Passage Retrieval (DPR) established
the now-standard retriever architecture underlying this and later work: a
pair of encoders, one for questions and one for passages, trained so that
relevant question-passage pairs are close in embedding space, which
outperformed classical sparse (term-matching) retrieval on open-domain QA
benchmarks when trained directly on question-passage relevance
\cite{karpukhin2020dpr}. Both lines of work depend on the Transformer
architecture's self-attention mechanism for the encoders and generators
involved \cite{vaswani2017attention}, and on sentence-embedding methods
such as Sentence-BERT that make large-scale semantic similarity search
computationally practical outside of a fine-tuned retrieval-specific
model \cite{reimers2019sbert}.

\subsection{Reader Architectures}
Given retrieved passages, how a generator consumes them matters as much
as how they were retrieved. Fusion-in-Decoder (FiD) showed that encoding
each retrieved passage independently and fusing evidence only in the
decoder's cross-attention lets a generative reader scale to many more
retrieved passages than concatenation-based encoding allows, without a
quadratic attention cost, and that answer accuracy continues to improve
as more passages are retrieved under this scheme \cite{izacard2021fid}.
Self-RAG moved a step further by training the model itself to decide,
via learned reflection tokens, when retrieval is necessary at all and to
critique its own output against retrieved evidence, rather than treating
retrieval as an unconditional, fixed pipeline stage that always runs
before generation \cite{asai2024selfrag}. Read together, these two works
mark a shift from ``retrieve a fixed $k$ passages, then generate'' toward
readers that are increasingly aware of how much and which evidence they
are actually using.

\subsection{Vector Search Infrastructure}
All of the above requires an index that can return the nearest neighbours
of a query embedding among potentially millions of candidate vectors
quickly enough for interactive use. FAISS demonstrated that approximate
nearest-neighbour search scales to billion-vector collections with
acceptable recall/latency trade-offs on commodity GPU hardware
\cite{johnson2017faiss}, and remains the substrate under many later,
higher-level vector stores. Developer-facing systems such as Chroma
build on the same underlying idea but expose it as a collection-oriented
API with metadata filtering, designed to be embedded directly inside an
application rather than operated as a standalone service
\cite{chroma2026docs}. This shift -- from ``a nearest-neighbour index'' to
``an application-embeddable vector store with metadata filtering'' -- is
one of the more consequential but under-discussed enablers of the recent
wave of RAG applications, since it is what lets an application filter
retrieval by, for instance, which knowledgebase or source a chunk came
from, rather than only by embedding similarity.

\subsection{Graph-Based and Global-Question RAG}
Plain chunk retrieval is well suited to locally-answerable questions but
poorly suited to \emph{global}, corpus-level questions such as ``what are
the main themes across this entire source'', since no single retrieved
chunk is likely to represent the whole. GraphRAG addresses this directly
by using an LLM at index time to derive an entity-relationship graph
from the source corpus and to pregenerate hierarchical community
summaries over that graph, which are then used to answer global,
query-focused summarisation questions that conventional RAG
under-serves \cite{edge2024graphrag}. This is, to our knowledge, the most
directly comparable prior approach to the broad-question problem that
motivates SAGE's deterministic overview chunks (Section IX): GraphRAG
solves the same problem with a substantially heavier index-time
pipeline -- an LLM call (or many) per corpus to build and summarise the
graph -- while SAGE trades some of GraphRAG's summarisation depth for a
non-generative, source-derived overview that adds no second
hallucination-prone stage at indexing time. We view the two as
occupying different points on the same cost/depth trade-off rather than
as competitors, and return to this contrast in Section IX.

\subsection{Survey-Level Views of the Field}
Broader surveys of RAG as a field consistently identify the same three
concerns as dominant in real deployments regardless of the specific
architecture used: provenance (can a user trace an answer back to its
evidence), updateability (can the corpus change without retraining the
model), and modularity (can retrieval, reranking, and generation be
swapped independently) \cite{gao2023ragsurvey}. More recent systematic
reviews reach similar conclusions while additionally emphasising
evaluation methodology as an open problem in its own right -- that is,
that the field has produced many architectures faster than it has
produced agreed-upon ways to measure whether one architecture is
actually better than another for a given task \cite{slr2025rag,ragevalsurvey2025}.
This survey adopts provenance, modularity, and (per Section VI)
faithfulness as the organising concerns for the modality-specific and
applied sections that follow.

\section{Modality-Specific Retrieval and Understanding}

A single-modality RAG system can retrieve fine-grained chunks and stop.
A multi-modal assistant must additionally decide, per modality, what the
unit of retrieval and citation actually is, since ``a paragraph'' is not
the natural evidence unit for a video or a repository.

\subsection{Documents}
For visually structured documents, LayoutLM showed that jointly modelling
text content and 2-D layout position improves document understanding
tasks relative to text-only models, since meaning in forms, tables, and
structured reports depends on spatial arrangement as much as on token
sequence \cite{xu2020layoutlm}. This line of work is most relevant to
scanned or visually rich PDFs; for born-digital PDFs with selectable
text, the layout signal is less critical and page-level text extraction
is sufficient for reliable retrieval, which is the regime most practical
document-QA systems -- and the SAGE case study in Section IX -- currently
target, deferring OCR and layout-aware extraction to future work.

\subsection{Video and Speech}
Whisper demonstrated that a single, large, weakly-supervised speech
model can achieve robust multilingual transcription without
task-specific fine-tuning, which made local, dependency-light
transcription practical for downstream applications for the first time
at that level of robustness \cite{radford2023whisper}. In practice,
systems that ingest video content for QA purposes tend to combine two
paths: retrieving existing captions when available, and falling back to
a local Whisper-family model only when captions do not exist, since
transcription is the more expensive and less accurate of the two options.
The retrieval unit for video is naturally a timestamped transcript span
rather than a page or paragraph, which has direct consequences for how a
citation must be rendered (Section IX).

\subsection{Code}
Code introduces retrieval considerations that do not apply to prose:
relevance depends on file structure, line locality, and cross-file
context, not only on token-level semantic similarity. CodeBERT
established that a single model can learn strong joint representations
of natural language and code well enough to support natural-language-to-code
retrieval tasks \cite{feng2020codebert}. RepoCoder went further, showing
that retrieving relevant context from \emph{across the repository} --
not just the current file -- substantially improves completion and
generation tasks, because a large share of the context a correct
completion needs (helper functions, type definitions, project
conventions) is not visible in the local file at all \cite{zhang2023repocoder}.
This repository-level-context finding is, in our reading, the single
most directly transferable idea from the code-intelligence literature to
a general-purpose multi-source assistant: it motivates generating an
explicit repository-level overview alongside file-level chunks, rather
than treating a repository as an undifferentiated bag of files.

\subsection{Websites}
Website ingestion is comparatively under-theorised in the academic
literature relative to documents, video, and code, in part because it is
treated as a solved crawling/scraping problem rather than a retrieval
problem. In practice, two considerations dominate: many modern pages
require client-side rendering to expose their real text content, which
argues for browser automation (e.g.\ Playwright \cite{playwright2026})
over raw HTTP fetching; and a crawler that follows links on a
user-supplied URL is, unless explicitly constrained, both a scope-creep
risk (crawling far beyond the page the user actually cared about) and a
security risk, discussed directly in Section VII.

\section{Multi-Source and Notebook-Style Assistants}

\subsection{Consumer Assistants}
Google's NotebookLM packages several of the above modalities -- documents,
web pages, and video/audio among them -- behind a single consumer-facing,
source-grounded assistant, explicitly marketed on the premise that
grounding responses in user-supplied sources reduces hallucination
relative to an open-domain chatbot \cite{notebooklm2026}. As a closed
product, its retrieval, reranking, and citation-rendering internals are
not published, which makes it useful as a reference point for what a
multi-source assistant can look like to an end user, but not directly
reusable as a technical baseline for a research-oriented system -- a gap
this survey's case study (Section IX) is positioned to partially fill by
describing a comparable, fully open system.

\subsection{Orchestration Frameworks}
Where NotebookLM is a fixed product, LangChain \cite{langchain2026},
LlamaIndex \cite{llamaindex2026}, and Haystack \cite{haystack2026} are
general-purpose libraries: they supply connectors for many data sources,
composable indexing and query abstractions, and pipeline components for
retrieval, reranking, and generation, from which a developer assembles a
specific application. This generality is their central strength and also
their central limitation for the purposes of this survey: none of the
three frameworks ships a fixed, opinionated answer to the four design
questions raised in Section I (broad-question retrievability,
per-modality citation format, multi-source fairness, and ingestion-time
security) -- each is left to the application developer to decide, and in
practice is frequently left undecided, which is itself a notable,
under-discussed gap between what these frameworks make \emph{possible}
and what they make \emph{the default}.

\section{Confidence, Faithfulness, and Hallucination in RAG}

RAG reduces, but does not eliminate, unsupported generation: a retriever
can return irrelevant or only partially relevant passages, and a
generator can still produce claims unsupported by whatever was retrieved.
A substantial literature studies this directly. Comprehensive surveys of
hallucination in large language models organise the failure modes into
taxonomies spanning input-conflicting, context-conflicting, and
fact-conflicting hallucination, and catalogue detection and mitigation
techniques operating at the data, training, and inference stages
respectively \cite{huang2024hallucination}. Detection techniques that are
particularly relevant to a retrieval-augmented setting include
sampling-based, black-box consistency checks that flag likely
hallucinations by testing whether repeated generations from the same
model remain internally consistent, without requiring access to model
internals or a labelled dataset -- a property that makes such techniques
attractive for a system, like most practical RAG deployments, built on a
hosted, closed LLM API rather than a locally fine-tuned model.

Confidence estimation in deployed RAG systems tends to be considerably
simpler than this literature: rather than a calibrated posterior
probability, most practical systems -- including the SAGE case study in
Section IX -- compute an interpretable heuristic from retrieval and
reranking signals (best score, evidence agreement, citation count) and
use it to gate whether the system answers at all, or to display a
confidence label to the end user. This is a defensible engineering
compromise given the cost of true calibration, but it also means the
confidence values surfaced by most deployed assistants should be read as
operational guidance rather than as a rigorous faithfulness guarantee --
a distinction we found was rarely made explicit in product-facing
documentation, though it is well understood in the research literature
\cite{huang2024hallucination}.

\section{Security Considerations in RAG Pipelines}

A RAG system that retrieves external content before generating a
response has an attack surface that a closed, purely parametric chatbot
does not, and this surface grows directly with the number of external
source types the system is willing to ingest.

\subsection{Indirect Prompt Injection}
Greshake et al.\ formalised \emph{indirect prompt injection}: an attacker
does not interact with the LLM-integrated application directly, but
instead places adversarial instructions inside content -- a web page, a
document, a retrieved search result -- that the application will later
fetch and place into the model's context, at which point the model may
treat those instructions as if they came from the application's own
system prompt or the user \cite{greshake2023injection}. The paper's
threat taxonomy (data exfiltration, worming across chained agents,
ecosystem contamination, and unauthorised tool or API invocation) applies
directly to any multi-source assistant, since every ingested website,
PDF, transcript, or repository file is, from the model's perspective,
content the system chose to trust enough to place in context. This risk
is now formally ranked as the top item in the OWASP Top 10 for LLM
Applications \cite{owasp2025llm01}, reflecting a broader industry
consensus that it is not a theoretical concern but the most commonly
exploited weakness in deployed LLM-integrated applications.

\subsection{Server-Side Request Forgery in Ingestion}
A second, less frequently discussed risk is specific to systems that
accept a user-supplied URL as an ingestion target: unless every request
the system makes on that URL's behalf -- including automatic HTTP
redirects and, for browser-automated crawlers, in-page navigations
triggered after the initial request -- is validated against private,
loopback, link-local, and cloud-metadata address ranges, the ingestion
endpoint becomes a general-purpose SSRF primitive that can be used to
probe or reach infrastructure the operator did not intend to expose. This
risk is well documented in general web-application security guidance but
is, in our reading of the RAG-specific literature, rarely discussed
alongside indirect prompt injection despite affecting the same class of
system for the same underlying reason: both stem from treating
ingestion-time content as fully trusted once a single, initial check has
passed.

\subsection{Implications for System Design}
Taken together, these two risks argue for security controls that are
structural rather than input-filtering-based: revalidating every network
hop rather than only the user-supplied URL, and delimiting retrieved
content from instructions at the prompt level rather than attempting to
detect and strip injected instructions after the fact, since detection-based
defenses against injected instructions are themselves an active and
unsettled research problem \cite{greshake2023injection,owasp2025llm01}.
Section IX describes how the SAGE case study implements both controls
concretely.

\section{Hybrid and Advanced Retrieval Strategies}

Dense, embedding-based retrieval and classical sparse, term-based
retrieval (e.g.\ BM25) fail on complementary classes of query: dense
retrieval generalises well across paraphrase and synonymy but can miss
passages that share no semantic neighbourhood with the query despite
containing an exact match for a code identifier, error string, or proper
noun, while sparse retrieval is precise on exact terms but blind to
paraphrase. A substantial and consistent empirical finding across the
retrieval literature is that combining the two -- for example via
Reciprocal Rank Fusion (RRF), which combines multiple rankings by summing
reciprocal-rank scores rather than requiring the rankings' underlying
scores to be comparable -- reliably outperforms either method used alone,
and does so without requiring a trained fusion model
\cite{cormack2009rrf}. This finding predates the current wave of
LLM-centric RAG systems by over a decade, but is, in our assessment,
under-adopted in practice: many deployed RAG pipelines, including
single-vector-store architectures like the SAGE case study discussed
next, currently rely on dense retrieval alone or on a fixed linear blend
of dense and lexical scores rather than a full hybrid pipeline, which we
identify in Section X as one of the more immediately actionable gaps
across the systems surveyed here.

\section{Case Study: SAGE}

We use SAGE, a prototype we built and validated, to make the design
space of Sections III--VIII concrete. SAGE ingests four source types --
public websites, born-digital PDFs, YouTube videos, and public GitHub
repositories -- into a shared chunk-and-metadata representation: one
knowledgebase per source, a deterministic overview chunk generated from
source statistics and excerpts (not a second LLM call) followed by
detailed chunks, and provenance metadata specific to the modality (URL
and chunk index for websites, page number for PDFs, timestamp span for
video, file path and line range for GitHub). Unlike GraphRAG's
LLM-built entity graph and community summaries \cite{edge2024graphrag},
SAGE's overview chunk is compiled deterministically from source
statistics and excerpts, trading summarisation depth for a
non-generative pipeline stage that cannot itself introduce new,
unsupported claims -- a design choice appropriate to a system whose
corpora are ingested from arbitrary, untrusted public sources rather
than a curated private dataset. Retrieval embeds chunks with
a compact sentence encoder \cite{minilm2026}, stores vectors in Chroma
\cite{chroma2026docs}, and stores relational metadata and lifecycle state
in Supabase PostgreSQL \cite{supabase2026docs}. At query time, SAGE
optionally rewrites short or ambiguous questions using recent
conversation turns, retrieves and reranks candidate chunks with a linear
blend of semantic and lexical scores, and -- for questions that combine
more than one selected source -- applies a diversity-floor step that
reserves each source's best-scoring chunk before filling remaining slots
by pure score, guaranteeing every selected source is represented in the
final citation set rather than only the highest-scoring one. Generation
runs in a Grounded Mode that refuses when retrieval evidence is weak, or
an Exploratory Mode that presents source-backed content before any
broader explanation; both modes wrap retrieved content in explicit,
delimited tags with an instruction-injection guard directly motivated by
\cite{greshake2023injection}, and ingestion validates every fetched URL
and redirect hop against private/loopback/link-local/cloud-metadata
address ranges.

Table~\ref{tab:axes} maps SAGE's concrete design choices onto the four
organising questions from Section I, alongside NotebookLM and a
representative application built on a general-purpose framework.

\begin{table}[t]
\caption{The four organising design questions, answered by SAGE, NotebookLM \cite{notebooklm2026}, and a representative application built on a general-purpose framework \cite{langchain2026,llamaindex2026,haystack2026}.}
\label{tab:axes}
\centering
\begin{tabular}{p{1.6cm}p{1.7cm}p{1.7cm}p{1.7cm}}
\hline
\textbf{Question} & \textbf{SAGE} & \textbf{NotebookLM} & \textbf{Framework app} \\
\hline
Broad-question retrievability & Deterministic overview chunk per source & Not disclosed & Left to developer \\
Per-modality citation & URL, page, timestamp, or file/line (fixed per modality) & Citations shown; internals undisclosed & Left to developer \\
Multi-source fairness & Diversity-floor rerank (Alg.\ guaranteed) & Not disclosed & Left to developer \\
Ingestion \& generation security & Explicit per-hop SSRF check + injection-guarded prompting & Not disclosed & Left to developer \\
\hline
\end{tabular}
\end{table}

The pattern in Table~\ref{tab:axes} is the survey's central observation
about the current state of the field: none of these four questions is
architecturally hard to answer well, but only a fixed, opinionated
system -- whether a closed consumer product or an open prototype like
SAGE -- actually answers all four by default, while a general-purpose
framework leaves all four to whoever assembles the final application,
which in our reading of deployed framework-based projects is
inconsistently done at best.

\section{Open Challenges and Future Directions}

Four gaps recur across the systems and literature surveyed here, and
remain open in the SAGE case study as much as elsewhere. First,
\emph{calibrated confidence}: nearly every practical system, SAGE
included, substitutes an interpretable heuristic for a calibrated
faithfulness estimate, and closing this gap requires either adopting
sampling-based hallucination-detection techniques from the literature
\cite{huang2024hallucination} at inference time, at some latency and cost
premium, or investing in a controlled, annotated evaluation set specific
to the deployment domain. Second, \emph{hybrid retrieval}: despite over a
decade of evidence that combining dense and sparse retrieval outperforms
either alone \cite{cormack2009rrf}, single-vector-store architectures
remain common, largely because they are simpler to operate; the actual
engineering cost of adding a sparse index alongside an existing dense
one is modest relative to the precision gained on exact-match queries.
Third, \emph{security as a first-class evaluation criterion}: indirect
prompt injection and ingestion-time SSRF are both well documented as
risks \cite{greshake2023injection,owasp2025llm01}, but we found no
widely-adopted benchmark that scores a RAG system's resistance to either,
which leaves builders to implement structural defenses (as SAGE does)
without a standard way to verify their sufficiency. Fourth,
\emph{cross-modal evaluation}: existing RAG benchmarks are overwhelmingly
single-modality (document QA or code completion, rarely both, and almost
never alongside video), which makes it difficult to compare multi-modal
assistants like SAGE and NotebookLM on any common, reproducible ground
rather than on vendor-reported or informally validated results.

\section{Conclusion}

This survey reviewed retrieval-augmented generation from its core
architectures through modality-specific extraction, multi-source
assistants and orchestration frameworks, confidence and hallucination,
and the security of the retrieval pipeline itself, and organised the
resulting design space around four questions: broad-question
retrievability, per-modality citation provenance, multi-source fairness,
and ingestion/generation security. Using SAGE as a case study showed
that all four questions admit concrete, implementable answers --
deterministic overview chunks, modality-aware citation metadata,
diversity-floor reranking, and combined SSRF and injection defenses --
at a cost that is modest relative to building any RAG pipeline at all,
which suggests the gap between a merely functional RAG system and a
trustworthy, source-grounded one is primarily a design-discipline
problem rather than a research problem requiring new architectures.
Calibrated confidence, hybrid retrieval, security benchmarking, and
cross-modal evaluation remain the open problems most likely to determine
whether the next generation of these systems closes that gap further.

\balance
\begin{thebibliography}{99}

\bibitem{lewis2020rag} P. Lewis \emph{et al.}, ``Retrieval-Augmented Generation for Knowledge-Intensive NLP Tasks,'' in \emph{Advances in Neural Information Processing Systems}, vol. 33, 2020, pp. 9459--9474.

\bibitem{karpukhin2020dpr} V. Karpukhin \emph{et al.}, ``Dense Passage Retrieval for Open-Domain Question Answering,'' in \emph{Proc. EMNLP}, 2020, pp. 6769--6781.

\bibitem{reimers2019sbert} N. Reimers and I. Gurevych, ``Sentence-BERT: Sentence Embeddings using Siamese BERT-Networks,'' in \emph{Proc. EMNLP-IJCNLP}, 2019, pp. 3982--3992.

\bibitem{vaswani2017attention} A. Vaswani \emph{et al.}, ``Attention Is All You Need,'' in \emph{Advances in Neural Information Processing Systems}, vol. 30, 2017, pp. 5998--6008.

\bibitem{izacard2021fid} G. Izacard and E. Grave, ``Leveraging Passage Retrieval with Generative Models for Open Domain Question Answering,'' in \emph{Proc. 16th Conf. European Chapter of the Association for Computational Linguistics (EACL)}, 2021, pp. 874--880.

\bibitem{asai2024selfrag} A. Asai, Z. Wu, Y. Wang, A. Sil, and H. Hajishirzi, ``Self-RAG: Learning to Retrieve, Generate, and Critique through Self-Reflection,'' in \emph{Proc. Int. Conf. on Learning Representations (ICLR)}, 2024.

\bibitem{edge2024graphrag} D. Edge \emph{et al.}, ``From Local to Global: A Graph RAG Approach to Query-Focused Summarization,'' \emph{arXiv preprint arXiv:2404.16130}, 2024.

\bibitem{johnson2017faiss} J. Johnson, M. Douze, and H. J\'egou, ``Billion-scale Similarity Search with GPUs,'' \emph{arXiv preprint arXiv:1702.08734}, 2017.

\bibitem{chroma2026docs} Chroma, ``Chroma Documentation,'' 2026. [Online]. Available: \url{https://docs.trychroma.com}. Accessed: Jul. 10, 2026.

\bibitem{gao2023ragsurvey} Y. Gao \emph{et al.}, ``Retrieval-Augmented Generation for Large Language Models: A Survey,'' \emph{arXiv preprint arXiv:2312.10997}, 2023.

\bibitem{slr2025rag} Anonymous, ``A Systematic Literature Review of Retrieval-Augmented Generation: Techniques, Metrics, and Challenges,'' \emph{arXiv preprint arXiv:2508.06401}, 2025.

\bibitem{ragevalsurvey2025} Anonymous, ``Retrieval Augmented Generation Evaluation in the Era of Large Language Models: A Comprehensive Survey,'' \emph{arXiv preprint arXiv:2504.14891}, 2025.

\bibitem{xu2020layoutlm} Y. Xu \emph{et al.}, ``LayoutLM: Pre-training of Text and Layout for Document Image Understanding,'' in \emph{Proc. ACM SIGKDD}, 2020, pp. 1192--1200.

\bibitem{radford2023whisper} A. Radford \emph{et al.}, ``Robust Speech Recognition via Large-Scale Weak Supervision,'' in \emph{Proc. ICML}, vol. 202, 2023, pp. 28492--28518.

\bibitem{feng2020codebert} Z. Feng \emph{et al.}, ``CodeBERT: A Pre-Trained Model for Programming and Natural Languages,'' in \emph{Findings of EMNLP}, 2020, pp. 1536--1547.

\bibitem{zhang2023repocoder} F. Zhang \emph{et al.}, ``RepoCoder: Repository-Level Code Completion Through Iterative Retrieval and Generation,'' in \emph{Proc. EMNLP}, 2023, pp. 2471--2484.

\bibitem{playwright2026} Microsoft, ``Playwright Documentation,'' 2026. [Online]. Available: \url{https://playwright.dev/}. Accessed: Jul. 10, 2026.

\bibitem{notebooklm2026} Google, ``NotebookLM: AI Research Tool and Thinking Partner,'' 2026. [Online]. Available: \url{https://notebooklm.google/}. Accessed: Jul. 10, 2026.

\bibitem{langchain2026} LangChain, ``LangChain Documentation,'' 2026. [Online]. Available: \url{https://www.langchain.com/}. Accessed: Jul. 10, 2026.

\bibitem{llamaindex2026} LlamaIndex, ``LlamaIndex Documentation,'' 2026. [Online]. Available: \url{https://docs.llamaindex.ai/}. Accessed: Jul. 10, 2026.

\bibitem{haystack2026} deepset, ``Haystack: Open-Source LLM Orchestration Framework,'' 2026. [Online]. Available: \url{https://github.com/deepset-ai/haystack}. Accessed: Jul. 10, 2026.

\bibitem{huang2024hallucination} L. Huang \emph{et al.}, ``A Survey on Hallucination in Large Language Models: Principles, Taxonomy, Challenges, and Open Questions,'' \emph{ACM Transactions on Information Systems}, 2024.

\bibitem{greshake2023injection} K. Greshake, S. Abdelnabi, S. Mishra, C. Endres, T. Holz, and M. Fritz, ``Not What You've Signed Up For: Compromising Real-World LLM-Integrated Applications with Indirect Prompt Injection,'' in \emph{Proc. 16th ACM Workshop on Artificial Intelligence and Security (AISec)}, 2023.

\bibitem{owasp2025llm01} OWASP Gen AI Security Project, ``LLM01:2025 Prompt Injection,'' OWASP Top 10 for LLM Applications, 2025. [Online]. Available: \url{https://genai.owasp.org/llmrisk/llm01-prompt-injection/}. Accessed: Jul. 10, 2026.

\bibitem{cormack2009rrf} G. V. Cormack, C. L. A. Clarke, and S. Buettcher, ``Reciprocal Rank Fusion Outperforms Condorcet and Individual Rank Learning Methods,'' in \emph{Proc. 32nd Int. ACM SIGIR Conf.}, 2009, pp. 758--759.

\bibitem{minilm2026} Sentence Transformers, ``all-MiniLM-L6-v2 Model Card,'' 2026. [Online]. Available: \url{https://huggingface.co/sentence-transformers/all-MiniLM-L6-v2}. Accessed: Jul. 10, 2026.

\bibitem{supabase2026docs} Supabase, ``Database Overview,'' 2026. [Online]. Available: \url{https://supabase.com/docs/guides/database/overview}. Accessed: Jul. 10, 2026.

\end{thebibliography}

\end{document}
